\documentclass{standalone}
\usepackage{tikz}
\usetikzlibrary{shapes.geometric, arrows, positioning}

\begin{document}

\tikzstyle{iniciofim} = [ellipse, draw, fill=gray!20, text width=3cm, text centered, minimum height=1cm]
\tikzstyle{processo} = [rectangle, draw, fill=blue!15, text width=4cm, text centered, rounded corners, minimum height=1cm]
\tikzstyle{decisao} = [diamond, draw, fill=green!10, text width=3.5cm, text centered, aspect=2]
\tikzstyle{seta} = [->, thick]

\begin{tikzpicture}[node distance=1.6cm]

\node (inicio) [iniciofim] {Início};
\node (leitura) [processo, below of=inicio] {Leitura dos sensores DHT11 e Higrômetro};
\node (validar) [decisao, below of=leitura, yshift=-0.5cm] {Leitura válida?};
\node (salvar) [processo, below of=validar, yshift=-0.5cm] {Salvar dados no cartão SD};
\node (loop) [processo, below of=salvar, yshift=-0.5cm] {Aguardar intervalo e repetir};
\node (fim) [iniciofim, below of=loop, yshift=-0.5cm] {Fim};

\draw [seta] (inicio) -- (leitura);
\draw [seta] (leitura) -- (validar);
\draw [seta] (validar) -- node[right]{Sim} (salvar);
\draw [seta] (salvar) -- (loop);
\draw [seta] (loop) -- ++(3,0) |- (leitura);
\draw [seta] (validar.west) -- ++(-2,0) node[left]{Não} |- (leitura.west);
\draw [seta] (loop) -- (fim);

\end{tikzpicture}

\end{document}

